%!TEX root = ../main.tex
\section{Data and Lexicons}
\label{sec:data}

\subsection{Vietnamese Corpus (ViHSD)}
We use UIT-ViHSD \cite{luu2021vihsd}, a corpus of Vietnamese social-media
comments annotated with three mutually exclusive labels: \clean{}, \off{}, and \hate{}.
After cleaning, the corpus contains 33{,}398 comments with fixed splits: train 24{,}046,
validation 2{,}672, test 6{,}680.  The test split is imbalanced: 83.05\% \clean{},
10.30\% \hate{}, 6.65\% \off{}.  Macro-F1 is primary metric; per-class F1 is always reported.

\subsection{English Corpus (Davidson et al.)}
For English we use the Hate Speech and Offensive Language dataset
\cite{davidson2017automated}.  The fixed split: train 19{,}816, validation 2{,}477,
test 2{,}478.  The English test is also imbalanced but with a different profile:
77.8\% \off{}, 5.8\% \hate{}, 16.4\% \clean{}.

\subsection{Mixed Evaluation Sample}
To evaluate the bilingual pipeline end-to-end we use the complete held-out test
splits: 6{,}680 Vietnamese comments from ViHSD test and 2{,}478 English comments
from Davidson test, yielding 9{,}158 mixed bilingual samples (Table~\ref{tab:mixed_composition}).
This full-split evaluation preserves the per-language class priors.

\begin{table}[H]
\centering
\caption{Mixed evaluation sample composition ($n=9{,}158$ full test splits)}
\label{tab:mixed_composition}
\begin{tabularx}{\textwidth}{l@{\extracolsep{\fill}}ccccr}
\toprule
Language & Total & \hate & \off & \clean & Source \\
\midrule
Vietnamese & 6{,}680 & 688 & 444 & 5{,}548 & ViHSD test \\
English    & 2{,}478 & 143 & 1{,}918 & 417 & Davidson test \\
\midrule
\textbf{Total} & \textbf{9{,}158} & \textbf{831} & \textbf{2{,}362} & \textbf{5{,}965} & \\
\bottomrule
\end{tabularx}
\end{table}


\subsection{Vietnamese Lexicon (668 entries)}
The Vietnamese rule engine uses a hate-speech lexicon of \textbf{668 entries},
mined from training portions of ViHSD, R-ViHSD \cite{bontoriTdtu}, and ViCM
\cite{nguyen2026vicm}.  Most entries (82.5\%) are multi-word constructions;
78.4\% are context-required.  The lexicon contains 594 \hate{} and 74 \off{}
terms, explaining why the rule engine is weak on \off{}.

\subsection{English Lexicon (818 entries)}
The English lexicon contains \textbf{818 entries} (86 \hate{}, 356 \off{}, 376 \clean{}),
extracted from Davidson training split.  The English lexicon has far fewer HATE
terms (86 vs. 594) because Davidson corpus has only 5.8\% HATE documents.  This
asymmetry directly impacts the hybrid's HATE F1 (0.453 EN vs 0.610 VI) and is
the highest-priority target for future lexicon expansion.
